\documentclass[10pt,a4paper]{amsart}
\usepackage[margin=19mm]{geometry}
\usepackage{amsmath,amssymb,mathtools}
\usepackage[scr=rsfs,cal=euler]{mathalfa}
\usepackage{xfrac}
\usepackage[unicode,hidelinks,bookmarks=false]{hyperref}
\newcommand{\ie}{i.e.\ }
\newcommand{\cf}{cf.\ }
\newcommand{\resp}{respectively\ }
\newcommand{\Subsection}{\subsection}
\providecommand{\nobreakdash}{\nobreak\hbox{-}}
\setlength{\emergencystretch}{2em}
\multlinegap0pt
\usepackage[shortlabels]{enumitem}
\usepackage{amssymb}
\usepackage{tikz}
\usepackage{xcolor}
\usepackage{float}
\newtheorem{definition}{Definition}
\newtheorem{lemma}{Lemma}
\title{Free subgroups in weighted Leavitt Path Algebras}
\author{Huynh Viet Khanh}
\date{Source: arXiv:2606.20288v1 (CC BY 4.0)}
\begin{document}
\maketitle

\noindent\emph{Source and licence.} Free subgroups in weighted Leavitt Path Algebras. Version arXiv:2606.20288v1 (CC BY 4.0), distributed by arXiv under the Creative Commons Attribution 4.0 International licence, http://creativecommons.org/licenses/by/4.0/, as recorded in the arXiv licence metadata for this exact version and shown on the abstract page https://arxiv.org/abs/2606.20288v1. Full licence notice: \texttt{licenses/arxiv-2606.20288-CC-BY-4.0.txt}.\par\smallskip

\noindent\textit{Editorial scope.} Counted unit: the article's lemma lem:prescribed-lift (source0.tex lines 444--446). The article's complete original proof of it is delivered with it (source0.tex lines 448--566, one \texttt{proof} environment of 119 lines). No mathematics is written, repaired or re-derived: the statement and the proof are the article's own lines, in the article's own notation, and no retained line is substituted or re-flowed.  Retained uncounted as the source-local context the proof needs: lines 92--99.  Nothing was omitted from any retained span, so the package carries no omission record.  No retained line contains a citation, so the wrapper carries no bibliography and no bibliography entry.  No retained line contains a figure environment, an image-inclusion command or drawing code, so no image is delivered and the package declares no asset dependency.  Editorial formatting only, in addition: an AMS \texttt{amsart} preamble that restates the article's own declarations and macros used by the retained lines, namely the environments \texttt{definition}, \texttt{lemma} on separate counters, together with the standard packages those lines need.  The retained environments print in this document as Definition 1, Lemma 1 in order of appearance, whereas the article prints the same items with the numbering of its own full document.  The complete native source of the article as distributed in the arXiv e-print for this exact version is bundled unchanged as \texttt{sources/203203/source0.tex}.
\begin{definition}\label{definition_VF}
A \textit{representation graph} of $(E,\omega)$ is a pair $(F,\phi)$, where $F=(F^0,F^1,s_F,r_F)$ is a directed graph and $\phi=(\phi^0,\phi^1):F\to\widehat E$ is a graph homomorphism, such that:
\begin{enumerate}[label=\textup{(\arabic*)}]
\item for any $v\in F^0$ and $1\leq i\leq\omega(\phi^0(v))$, there is precisely one $f\in s_F^{-1}(v)$ with $\operatorname{tag}(\phi^1(f))=i$;
\item for any $v\in F^0$ and any $\mathtt{e}\in E^{\mathrm{st}}$ with $r(\mathtt{e})=\phi^0(v)$, there is precisely one $f\in r_F^{-1}(v)$ with $\operatorname{st}(\phi^1(f))=\mathtt{e}$.
\end{enumerate}
We refer these two conditions as \textit{outgoing condition} and \textit{incoming condition} respectively.
\end{definition}

\begin{lemma}\label{lem:prescribed-lift}
Let $(E,\omega)$ be a weighted graph, and let $\widehat E$ be the directed graph associated to $(E,\omega)$. Let $\mathtt{e}_i\in \widehat E^1$, where $\mathtt{e}\in E^{\mathrm{st}}$ and $1\leq i\leq \omega(\mathtt{e})$. Then there is a connected representation graph $(F,\phi)$ of $(E,\omega)$ containing an edge $e_i$ such that $\phi^1(e_i)=\mathtt{e}_i$.
\end{lemma}

\begin{proof}
Put $\mathtt{u}=s(\mathtt{e})$ and $\mathtt{v}=r(\mathtt{e})$. We begin with the smallest possible lift of the edge $\mathtt{e}_i$. Let $F_0$ be the directed graph with two vertices $u,v$ and one edge $e_i$ from $u$ to $v$. Define
$$
\phi_0^0(u)=\mathtt{u},\qquad
\phi_0^0(v)=\mathtt{v},\quad\text{and}\quad
\phi_0^1(e_i)=\mathtt{e}_i.
$$
Then $\phi_0:F_0\to \widehat E$ is a graph homomorphism.

The graph $F_0$ need not satisfy the two conditions in Definition~\ref{definition_VF}. We will enlarge it step by step. At each stage we add exactly the missing edges required by the outgoing and incoming conditions, while keeping uniqueness. That means we construct an increasing sequence of graphs inductively:
$$
(F_0,\phi_0)\subseteq (F_1,\phi_1)\subseteq (F_2,\phi_2)\subseteq\cdots .
$$

Assume that $(F_n,\phi_n)$ has been constructed. For $z\in F_n^0$ and $j\in\mathbb N$, define
$$
O_n(z,j)=
\{h\in F_n^1\mid s_{F_n}(h)=z,\ 
\operatorname{tag}(\phi_n^1(h))=j\}.
$$
Thus $O_n(z,j)$ records the edges in $F_n$ which start at $z$ and have tag $j$. Similarly, for $z\in F_n^0$ and $\mathtt{f}\in E^{\mathrm{st}}$ with $r(\mathtt{f})=\phi_n^0(z)$, define
$$
I_n(z,\mathtt{f})=
\{h\in F_n^1\mid r_{F_n}(h)=z,\ 
\operatorname{st}(\phi_n^1(h))=\mathtt{f}\}.
$$
Thus $I_n(z,\mathtt{f})$ records the edges in $F_n$ which end at $z$ and have structure edge $\mathtt{f}$. We shall make sure that, for every $n$,
\begin{equation}\label{eq:In-On}
|O_n(z,j)|\leq 1\qquad \text{and}\qquad |I_n(z,\mathtt{f})|\leq 1.
\end{equation}
These inequalities hold for $n=0$.

We now identify the defects at stage $n$. The outgoing defects are
$$
D_n^+=\{(z,j)\mid z\in F_n^0,\ 1\leq j\leq \omega(\phi_n^0(z)),
\ O_n(z,j)=\emptyset\}.
$$
The incoming defects are
$$
D_n^-=\{(z,\mathtt{f})\mid z\in F_n^0,\ \mathtt{f}\in E^{\mathrm{st}},
\ r(\mathtt{f})=\phi_n^0(z),\ I_n(z,\mathtt{f})=\emptyset\}.
$$
So $D_n^+$ lists the missing outgoing tags, while $D_n^-$ lists the missing incoming structure edges.

We first repair the outgoing defects. Let $(z,j)\in D_n^+$. Since $1\leq j\leq \omega(\phi_n^0(z))$, there exists a structured edge $\mathtt{f}(z,j)\in E^{\mathrm{st}}$ such that
$$
s(\mathtt{f}(z,j))=\phi_n^0(z),\quad\text{where}\quad
\omega(\mathtt{f}(z,j))\geq j.
$$
Add a new vertex $z_j^+$ and a new edge $f_{z,j}^+$ from $z$ to $z_j^+$. Define
$$
\phi_{n+1}^0(z_j^+)=r(\mathtt{f}(z,j))\quad\text{and}\quad
\phi_{n+1}^1(f_{z,j}^+)=\mathtt{f}(z,j)_j.
$$

We next repair the incoming defects. Let $(z,\mathtt{f})\in D_n^-$. Add a new vertex $z_{\mathtt{f}}^-$ and a new edge $f_{z,\mathtt{f}}^-$ from $z_{\mathtt{f}}^-$ to $z$. Define
$$
\phi_{n+1}^0(z_{\mathtt{f}}^-)=s(\mathtt{f})\quad\text{and}\quad
\phi_{n+1}^1(f_{z,\mathtt{f}}^-)=\mathtt{f}_1.
$$
All new vertices and new edges are taken to be distinct. On the old graph $F_n$, set $\phi_{n+1}=\phi_n$.

We now check that $\phi_{n+1}:F_{n+1}\to\widehat E$ is a graph homomorphism. For an edge $f_{z,j}^+$ added to repair an outgoing defect, we have
$$
\widehat s(\mathtt{f}(z,j)_j)=s(\mathtt{f}(z,j))=\phi_n^0(z)
$$
and
$$
\widehat r(\mathtt{f}(z,j)_j)=r(\mathtt{f}(z,j))=\phi_{n+1}^0(z_j^+).
$$
For an edge $f_{z,\mathtt{f}}^-$ added to repair an incoming defect, we have
$$
\widehat s(\mathtt{f}_1)=s(\mathtt{f})=\phi_{n+1}^0(z_{\mathtt{f}}^-)
$$
and
$$
\widehat r(\mathtt{f}_1)=r(\mathtt{f})=\phi_n^0(z).
$$
Thus the source and range are respected for every new edge, and hence $\phi_{n+1}$ is a graph homomorphism.

We now verify that the inequalities \eqref{eq:In-On} are preserved. Let $z\in F_n^0$ be an old vertex. For each missing outgoing tag $j$, we add exactly one edge starting at $z$ with tag $j$, and we add no edge for a tag already represented in $O_n(z,j)$. Hence $|O_{n+1}(z,j)|\leq 1$. Similarly, for each missing incoming structure edge $\mathtt{f}$, we add exactly one edge ending at $z$ with structure edge $\mathtt{f}$, and we add no edge for a structure edge already represented in $I_n(z,\mathtt{f})$. Hence $|I_{n+1}(z,\mathtt{f})|\leq 1$.

It remains to check the new vertices. First suppose that $z'=z_j^+$. Then $\phi_{n+1}^0(z')=r(\mathtt{f}(z,j))$, and the only edge of $F_{n+1}$ ending at $z'$ is $f_{z,j}^+$. Hence, for any $\mathtt{g}\in E^{\mathrm{st}}$ with $r(\mathtt{g})=\phi_{n+1}^0(z')$,
$$
I_{n+1}(z',\mathtt{g})=
\begin{cases}
\{f_{z,j}^+\},&\text{if }\mathtt{g}=\mathtt{f}(z,j),\\
\emptyset,&\text{if }\mathtt{g}\neq \mathtt{f}(z,j).
\end{cases}
$$
No edge of $F_{n+1}$ starts at $z'$, so
$$
O_{n+1}(z',\ell)=\emptyset
$$
for $1\leq \ell\leq \omega(\phi_{n+1}^0(z'))$. Thus the inequalities \eqref{eq:In-On} hold at $z'$.

Now suppose that $z'=z_{\mathtt{f}}^-$. Then $\phi_{n+1}^0(z')=s(\mathtt{f})$, and the only edge of $F_{n+1}$ starting at $z'$ is $f_{z,\mathtt{f}}^-$. Its image is $\mathtt{f}_1$. Hence, for $1\leq \ell\leq \omega(\phi_{n+1}^0(z'))$,
$$
O_{n+1}(z',\ell)=
\begin{cases}
\{f_{z,\mathtt{f}}^-\},&\text{if }\ell=1,\\
\emptyset,&\text{if }\ell\neq 1.
\end{cases}
$$
No edge of $F_{n+1}$ ends at $z'$, so $I_{n+1}(z',\mathtt{g})=\emptyset$ for every $\mathtt{g}\in E^{\mathrm{st}}$ with $r(\mathtt{g})=\phi_{n+1}^0(z')$. Thus the inequalities \eqref{eq:In-On} also hold at $z'$.

This completes the induction. Put
$$
F=\bigcup_{n\geq 0}F_n\qquad\text{and}\qquad
\phi=\bigcup_{n\geq 0}\phi_n.
$$
Since every new vertex is attached by an edge to a vertex already present, the graph $F$ is connected. The original edge $e_i$ from $u$ to $v$ remains in $F$, and $\phi^1(e_i)=\mathtt{e}_i$.

It remains only to check that $(F,\phi)$ is a representation graph. Let $z\in F^0$, and choose $n$ such that $z\in F_n^0$. If $1\leq j\leq \omega(\phi^0(z))$, then either $O_n(z,j)\neq\emptyset$, or $(z,j)\in D_n^+$. In the second case, the edge $f_{z,j}^+$ is added in $F_{n+1}$. Thus there is at least one edge $h\in s_F^{-1}(z)$ with $\operatorname{tag}(\phi^1(h))=j$. There is at most one such edge because $|O_m(z,j)|\leq 1$ for every finite $m$, and any two edges of $F$ already occur together in some $F_m$.

Similarly, let $\mathtt{f}\in E^{\mathrm{st}}$ satisfy $r(\mathtt{f})=\phi^0(z)$. Then either $I_n(z,\mathtt{f})\neq\emptyset$, or $(z,\mathtt{f})\in D_n^-$. In the second case, the edge $f_{z,\mathtt{f}}^-$ is added in $F_{n+1}$. Thus there is at least one edge $h\in r_F^{-1}(z)$ with $\operatorname{st}(\phi^1(h))=\mathtt{f}$. There is at most one such edge by the same finite-stage argument.

Therefore $(F,\phi)$ satisfies both the outgoing and incoming conditions in Definition~\ref{definition_VF}. Hence $(F,\phi)$ is a connected representation graph of $(E,\omega)$ containing the prescribed lift $e_i$ of $\mathtt{e}_i$.
\end{proof}

\end{document}
